\documentclass[10pt,a4paper]{amsart}
\usepackage[margin=19mm]{geometry}
\usepackage{amsmath,amssymb,mathtools}
\usepackage[scr=rsfs,cal=euler]{mathalfa}
\usepackage{xfrac}
\usepackage[unicode,hidelinks,bookmarks=false]{hyperref}
\newcommand{\ie}{i.e.\ }
\newcommand{\cf}{cf.\ }
\newcommand{\resp}{respectively\ }
\newcommand{\Subsection}{\subsection}
\providecommand{\nobreakdash}{\nobreak\hbox{-}}
\setlength{\emergencystretch}{2em}
\multlinegap0pt
\usepackage{amssymb}
\usepackage{xcolor}
\usepackage[normalem]{ulem}
\newtheorem{lemma}{Lemma}
\newtheorem{proposition}{Proposition}
\newcommand{\bZ}{\mathbb Z}
\def\ep{{\varepsilon}}
\title{Zero distribution of power series and binary correlation of coefficients}
\author{Jacques Benatar, Alexander Borichev, Mikhail Sodin}
\date{Source: arXiv:2104.04812v1 (CC BY 4.0)}
\begin{document}
\maketitle

\noindent\emph{Source and licence.} Zero distribution of power series and binary correlation of coefficients. Version arXiv:2104.04812v1 (CC BY 4.0), distributed by arXiv under the Creative Commons Attribution 4.0 International licence, http://creativecommons.org/licenses/by/4.0/, as recorded in the arXiv licence metadata for this exact version and shown on the abstract page https://arxiv.org/abs/2104.04812v1. Full licence notice: \texttt{licenses/arxiv-2104.04812-CC-BY-4.0.txt}.\par\smallskip

\noindent\textit{Editorial scope.} Counted unit: the article's proposition Lemma5a (source0.tex lines 2427--2452). The article's complete original proof of it is delivered with it (source0.tex lines 2454--2621, one \texttt{proof} environment of 168 lines). No mathematics is written, repaired or re-derived: the statement and the proof are the article's own lines, in the article's own notation, and no retained line is substituted or re-flowed.  Retained uncounted as the source-local context the proof needs: lines 1569--1577; lines 1588--1595; lines 1625--1637; lines 1669--1672; lines 1743--1747; lines 1768--1772; lines 1847--1870; lines 2392--2395.  Nothing was omitted from any retained span, so the package carries no omission record.  No retained line contains a citation, so the wrapper carries no bibliography and no bibliography entry.  No retained line contains a figure environment, an image-inclusion command or drawing code, so no image is delivered and the package declares no asset dependency.  Editorial formatting only, in addition: an AMS \texttt{amsart} preamble that restates the article's own declarations and macros used by the retained lines, namely the environments \texttt{lemma}, \texttt{proposition} on separate counters, together with the standard packages those lines need.  The retained environments print in this document as Lemma 1, Proposition 1 in order of appearance, whereas the article prints the same items with the numbering of its own full document.  The complete native source of the article as distributed in the arXiv e-print for this exact version is bundled unchanged as \texttt{sources/110147/source0.tex}.
\begin{lemma}\label{Lemma5A}
We have
\begin{itemize}
\item[{\rm (i)}]
$\sigma (s) = (1+o(1))\sigma $ everywhere on $[R, R(1+\beta)]$;
\item[{\rm (ii)}]
$\nu(R(1+\beta))-\nu (R) = (1+o(1))\beta\sigma$.
\end{itemize}
\end{lemma}

\begin{lemma}\label{Lemma5C}
We have
\begin{itemize}
\item[{\rm (i)}] $V_R(k; h)\lesssim \sqrt{\sigma}$;
\item[{\rm (ii)}] $V_R(k; 0) \gtrsim \sqrt{\sigma}$,
provided that  $\nu(R)\le k \le \nu(R(1+\beta))$.
\end{itemize}
\end{lemma}

\begin{lemma}\label{Lemma5B}\mbox{}
\begin{itemize}
\item[{\rm (i)}] $V_R(k; h) \lesssim e^{-ch^2/\sigma}\sqrt{\sigma}$;
\item[{\rm (ii)}]
\[
V_R(k; h) \lesssim
\begin{cases}
e^{-c(\nu(R)-k)^2/\sigma}\sqrt{\sigma}, \ &k<\nu(R), \\
e^{-c(k-\nu(R(1+\beta)))^2/\sigma}\sqrt{\sigma}, \ &k>\nu(R(1+\beta))\,.
\end{cases}
\]
\end{itemize}
\end{lemma}

\begin{lemma}\label{Lemma5D}
$ \displaystyle \sum_{k\in\bZ}
|V_R(k; h) - V_R(k-1; h)| \lesssim \sqrt{\sigma} $.
\end{lemma}

\begin{lemma}\label{Lemma10a} We have
\[
V_R(0) \simeq \beta \sigma^{3/2}\,.
\]
\end{lemma}

\begin{lemma}\label{Lemma10b} For $ h^2 \lesssim \sigma $, we have
\[
|V_R(h) - V_R(0)| \lesssim (1+ h^2\beta)\sqrt{\sigma}\,.
\]
\end{lemma}

\begin{proposition}\label{Lemma5}
Suppose that
\begin{equation}
\label{eq:condition1}
\sum_{\nu \le k \le \nu+\frac12 \beta\sigma}\, |\xi(k)|^2 \gtrsim \beta\sigma\,,
\end{equation}
and that, for some $p>1$,
\begin{equation}
\label{eq:condition2}
\sum_{h \ge 1} (1+\beta h)^{-p} S^*(M_1, M_2; h) \ll \beta\sigma,
\end{equation}
with
\[
M_1 = \nu (R) - \beta\sigma, \quad M_2 =  \nu(R(1+\beta)) + \beta\sigma\,.
\]
Then, for every $\vartheta\in [-\frac12, \frac12]$, there exist
\[
R'\in [R, R(1+\beta)], \quad \theta'\in (\vartheta-\beta, \vartheta+\beta),
\]
such that
\[
|W_{R'}(\theta')| \gtrsim \sigma^{1/4}\,.
\]
\end{proposition}

\begin{equation}\label{eq:C1}
\Bigl|\, \frac1X\, \sum_{0\le s < X} \xi(s)\bar\xi(s+h) - \widehat{\chi}(h) \Bigr|
\lesssim \ep_1(X; h), \qquad 0\le h \le H=H(X),
\end{equation}

\begin{proposition}\label{Lemma5a}
Let $R\gg 1$. Suppose that there exist $q>1$, $\beta=\beta(R)\ll 1/\Delta(\sigma)$,
and $g\in \mathcal G$, satisfying the following set of conditions:

\smallskip\noindent{\rm (a)}
$\beta^{-q} \le \min\bigl\{\sqrt{\sigma}, H(\tfrac12\, \nu)\bigr\}$;

\medskip\noindent{\rm (b)}
$\beta^{-(1+2q)} \ll \sigma\ep_2(\beta)$;

\medskip\noindent{\rm (c)}
$\displaystyle \nu \sum_{0\le h \le \beta^{-q}} \ep_1(\tfrac12\, \nu; h) \ll \beta\sigma\ep_2(\beta) $;

\medskip\noindent{\rm (d)}
$\displaystyle \sum_{h>\beta^{-q}} |\widehat{g}(\beta h)| \ll \ep_2(\beta) $.

\smallskip
Then, for every $\vartheta\in \bigl[-\tfrac12,  \tfrac12\, \bigr]$, there exist
\[
R'\in [R, R(1+\beta)],  \quad \theta'\in (\vartheta-\beta, \vartheta+\beta),
\]
such that
\[
|W_{R'}(\theta')| \gtrsim \sigma^{1/4}\, \sqrt{\ep_2 (\beta)}\,.
\]
\end{proposition}

\begin{proof}
As in the proof of Proposition~\ref{Lemma5}, we estimate from below the average
\[
\widetilde{X} =
\int_R^{R(1+\beta)}\, \int_{-1/2}^{1/2} |\widetilde{W}_s(\theta)|^2\,
g(\beta^{-1}(\vartheta-\theta))\, {\rm d}\theta\, {\rm d}\nu (s)\,,
\]
where
\[
\widetilde{W}_R(\theta) = \sum_{k\in\bZ} \xi(k)e(k\theta)
\exp\bigl[\,\frac{(k-\nu)^2}{2\sigma} \,\bigr]\,.
\]
By Lemma~\ref{Lemma5A}(ii), $\nu(R(1+\beta))-\nu(R) = (1+o(1))\beta\sigma$, so, to prove Proposition~\ref{Lemma5a} we need to show that
$\widetilde{X} \gtrsim \beta^2\sigma^{3/2}\ep_2(\beta)$.
As before, we rewrite $\widetilde X$ as a Fourier series
\[
\widetilde{X}
= \beta\,
\sum_{h\in\bZ} \widehat{g}(\beta h) e(-h\vartheta)\,
\sum_{k\in\bZ} \xi(k)\bar\xi(k+h) V_R(k; h)\,.
\]
Recalling the notation
$\displaystyle V_R(h) = \sum_{k\in\bZ} V_R(k; h)$, we split the RHS into three parts:
\begin{align*}
\widetilde{X}  &= \beta V_R(0)\, \sum_{h\in\bZ} \widehat{g}(\beta h)\widehat{\chi}(h) e(-h\vartheta) \\
&\quad + \beta\, \sum_{h\in\bZ} \widehat{g}(\beta h) e(-h\vartheta)\,
\sum_{k\in\bZ} \bigl( \xi(k)\bar\xi(k+h) - \widehat{\chi}(h) \bigr) V_R(k; h) \\
&\quad + \beta\, \sum_{h\in\bZ} \widehat{g}(\beta h) \widehat{\chi}(h) e(-h\vartheta)\,
(V_R(h) - V_R(0)) \\
&= I +II + III\,.
\end{align*}
We will show that $I\gtrsim \beta^2\sigma^{3/2}\ep_2(\beta)$, while
the terms $|II|$ and $|III|$ are  $\ll \beta^2\sigma^{3/2}\ep_2(\beta)$.

\medskip\noindent\underline{Lower bound of $I$}:
We set $g_\beta(\theta) = \beta^{-1}g(\beta^{-1}\theta)$, and denote by
$(\chi * g_\beta)'$ the density of the convolution of $\chi$ with
$g_\beta$. Then, \[ I=\beta V_R(0) (\chi * g_\beta)'(-\vartheta).\]
By Lemma~\ref{Lemma10a}, $V_R(0)\gtrsim \beta \sigma^{3/2}$. Since $g_\beta = \beta^{-1}$ on
$ \bigl[-\tfrac14\, \beta, \tfrac14\, \beta \bigr]$, we have
\[
(\chi * g_\beta)'(-\vartheta) \gtrsim
\chi\bigl[-\vartheta-\tfrac14\, \beta, -\vartheta + \tfrac14\, \beta \bigr]
\gtrsim \ep_2(\beta).
\]
Thus, $I \gtrsim \beta^2 \sigma^{3/2} \ep_2(\beta)$.

\medskip\noindent\underline{Upper bound of $II$}:
Recalling that $\beta\sigma \stackrel{(a)}\gg \sqrt{\sigma\log\sigma}$
and using Lemma~\ref{Lemma5B}(ii), we cut the sum in $k$, getting
\begin{align*}
|II| &\lesssim \beta\, \sum_{h\in\bZ} |\widehat{g}(\beta h)| \cdot
\bigl| \sum_{|k-\nu|\le 2\beta\sigma}
\bigl( \xi(k)\bar\xi(k+h) - \widehat{\chi}(h) \bigr) V_R(k; h)\, \bigr| + o(1) \\
&= \beta\, \bigl(\, \sum_{|h|\le \beta^{-q}} + \sum_{|h|>\beta^{-q}} \,\bigr)\,
|\widehat{g}(\beta h)| \\ &\qquad\qquad\qquad\qquad \times
\bigl|\, \sum_{|k-\nu|\le 2\beta\sigma}
\bigl( \xi(k)\bar\xi(k+h) - \widehat{\chi}(h) \bigr) V_R(k; h) \,\bigr| + o(1)\,.
\end{align*}
For $|h|>\beta^{-q}$, using the crude estimate
\[
\bigl|\, \sum_{|k-\nu|\le 2\beta\sigma}
\bigl( \xi(k)\bar\xi(k+h) - \widehat{\chi}(h) \bigr) V_R(k; h) \,\bigr|
\lesssim \beta\sigma \cdot \sqrt{\sigma} = \beta\sigma^{3/2}\,,
\]
we get
\begin{multline*}
\beta\, \sum_{|h|> \beta^{-q}}
|\widehat{g}(\beta h)| \cdot
\bigl|\, \sum_{|k-\nu|\le 2\beta\sigma}
\bigl( \xi(k)\bar\xi(k+h) - \widehat{\chi}(h) \bigr) V_R(k; h) \,\bigr|
\\
\lesssim \beta^2\sigma^{3/2} \sum_{|h|> \beta^{-q}}
|\widehat{g}(\beta h)|\,,
\end{multline*}
which is $\ll \beta^2\sigma^{3/2}\ep_2(\beta)$ by assumption (d).

Now, we consider the sum over $|h|\le\beta^{-q}$.
Applying summation by parts and using Lemma~\ref{Lemma5D}, we have
\begin{align*}
\Bigl| &\sum_{|k-\nu|\le 2\beta\sigma}
\bigl( \xi(k)\bar\xi(k+h) - \widehat{\chi}(h) \bigr) V_R(k; h) \Bigr|
\\
&\lesssim
\Bigl(\max_{|k-\nu|\le 2\beta\sigma} V_R(k; h)+ \sum_{|k-\nu|\le 2\beta\sigma} |V_R(k; h)-V_R(k-1; h)|\,\Bigr) 
\\&\qquad\qquad\qquad\qquad\qquad\times
\max_{|k-\nu|\le 2\beta}\, \bigl|\,
\sum_{k\le s \le \nu +2\beta\sigma}
\bigl(  \xi(s)\bar\xi(s+h) - \widehat{\chi}(h) \bigr)\,
\bigr|
\\
&\lesssim
\sqrt{\sigma}\, \max_{|k-\nu|\le 2\beta}\, \bigl|\,
\sum_{k\le s \le \nu +2\beta\sigma}
\bigl(  \xi(s)\bar\xi(s+h) - \widehat{\chi}(h) \bigr)
\,\bigr|\,.
\end{align*}
First, we assume that $0\le h \le \beta^{-q}$. Then, by estimate~\eqref{eq:C1},
the maximum on the RHS is $ \lesssim (\nu + 2\beta\sigma)\ep_1(\nu-2\beta\sigma; h)$.
Since $\beta\sigma \ll \sigma/\Delta(\sigma) \lesssim \nu$, the latter expression is
$ \lesssim \, \nu\ep_1(\tfrac12\, \nu; h) $. The application of estimate~\eqref{eq:C1} was legal since,
$  \beta^{-q} \stackrel{(a)}\le H(\tfrac12\, \nu) \,
\stackrel{\beta\sigma\ll\nu}\le\, H(\nu-2\beta\sigma) $.
Thus,
\begin{multline*}
\beta\, \sum_{0\le h \le \beta^{-q}}\,
|\widehat{g}(\beta h)| \cdot
\bigl|\, \sum_{|k-\nu|\le 2\beta\sigma}
\bigl( \xi(k)\bar\xi(k+h) - \widehat{\chi}(h) \bigr) V_R(k; h) \,\bigr|
\\
\lesssim
\beta\sqrt{\sigma}\cdot \nu\, \sum_{0\le h \le \beta^{-q}} \ep_1(\tfrac12\, \nu, h)\,,
\end{multline*}
which is $\ll \beta^2\sigma^{3/2} \ep_2(\beta)$ by assumption (c).

The sum over $-\beta^{-q} \le h \le -1$ needs only a minor modification.
In this case we have
\begin{align*}
\bigl|\, \sum_{k\le s \le \nu +2\beta\sigma}
\bigl(  \xi(s)\bar\xi(s+h) &- \widehat{\chi}(h) \bigr) \,\bigr|\\
&= \bigl|\, \sum_{k-|h|\le s \le \nu +2\beta\sigma-|h|}\,
\bigl(  \xi(s)\bar\xi(s+|h|) - \widehat{\chi}(|h|) \bigr) \,\bigr| \\
&= \bigl|\, \sum_{k\le s \le \nu +2\beta\sigma}\,
\bigl(  \xi(s)\bar\xi(s+|h|) - \widehat{\chi}(|h|) \bigr) \,\bigr| + O(|h|)\,.
\end{align*}
Therefore,
\begin{multline*}
\beta\, \sum_{-\beta^{-q}\le h \le -1}\,
|\widehat{g}(\beta h)| \cdot
\bigl|\, \sum_{|k-\nu|\le 2\beta\sigma}
\bigl( \xi(k)\bar\xi(k+h) - \widehat{\chi}(h) \bigr) V_R(k; h) \,\bigr|
\\
\lesssim
\beta\sqrt{\sigma}\bigl(\, \nu\, \sum_{0\le h' \le \beta^{-q}} \ep_1(\tfrac12\, \nu, h')
+ \beta^{-2q}\,\bigr)\,,
\end{multline*}
and, by assumptions (c) and (b), both terms on the RHS are $\ll \beta^2\sigma^{3/2} \ep_2(\beta)$.

\medskip\noindent\underline{Upper bound of $III$}:
We have
\[
| III | \lesssim \beta\, \bigl(\, \sum_{|h|\le\sqrt\sigma} + \sum_{|h|>\sqrt{\sigma}} \,\bigr)
|\widehat{g}(\beta h)|\cdot |(V_R(h) - V_R(0))|\,.
\]
To estimate the first sum, we apply Lemma~\ref{Lemma10b} and use that the Fourier transform of $g$ decays faster than any negative power. We get
\begin{multline*}
\beta\, \sum_{|h|\le\sqrt\sigma}
|\widehat{g}(\beta h)|\cdot |(V_R(h) - V_R(0))| \lesssim \beta\sqrt{\sigma}\,
\sum_{|h|\le\sqrt\sigma}\,  |\widehat{g}(\beta h)| (1+ h^2\beta) \\
\lesssim \beta\sqrt{\sigma} \cdot \beta^{-2} = \frac{\sqrt{\sigma}}\beta\,.
\end{multline*}
To estimate the second sum, we use the crude bound $V_R(h) \lesssim \beta\sigma^{3/2}$,
which follows from Lemma~\ref{Lemma5C}(i) and Lemma~\ref{Lemma5B}(ii).
Using again that the Fourier transform of $g$ decays faster than any negative power, we get
\begin{multline*}
\beta\, \sum_{|h|>\sqrt\sigma}
|\widehat{g}(\beta h)|\cdot |(V_R(h) - V_R(0))| \lesssim \beta^2\sigma^{3/2}\,
\sum_{|h|>\sqrt\sigma} |\widehat{g}(\beta h)| \\
\lesssim \beta^2 \sigma^{3/2} \cdot \frac1{\beta}\,
\sum_{\ell\ge\beta\sqrt\sigma} \frac1{\ell^3} \lesssim \frac{\sqrt{\sigma}}\beta\,.
\end{multline*}
It remains to recall that
condition (b) guarantees that 
$$
\beta^{-1}\sqrt{\sigma} \ll \beta^2 \sigma^{3/2}\, \ep_2(\beta).
$$
This completes the proof of Proposition~\ref{Lemma5a}.
\end{proof}

\end{document}
